\item \emph{Integral chains and a homotopy in every degree.}
Write $L_k=\bigwedge^k\mathbb Z^6$ with zero differential and
$W_{i,k}=\bigwedge^kM_i$. Let $\mathcal P_*(T^6)$ be integral
rational polyhedral chains, modulo oriented subdivision. Cut a
rational simplex at the integer coordinate hyperplanes, project to
the fundamental cube, and identify opposite faces. A finite set of
these chains has a common rational subdivision. Dimension-dropping
affine images are zero in that dimension. The resulting complex
computes integral torus homology by the subdivision equivalences of
finite triangulations. It is free in every degree: choose a countable
cofinal sequence of rational subdivisions; each subdivision inclusion
splits on chain groups by selecting one subdividing simplex in the
interior of each old simplex and recovering its coefficient there.

Addition on the torus gives the Pontryagin product
$u*v=\operatorname{add}_*(u\times v)$, with
\[
 \partial(u*v)=\partial u*v+(-1)^{\deg u}u*\partial v.
\]
Let $\gamma_j$ be the coordinate circle and $\ell_v$ the straight
lattice loop $t\mapsto tv$. The integral cycle inclusion is
\[
 \iota(e_{i_1}\wedge\cdots\wedge e_{i_k})
       =\gamma_{i_1}*\cdots*\gamma_{i_k},\qquad \iota(1)=[0].
\]
It induces the product-cell homology basis. Repeated coordinate
factors have dimension-dropping image, and interchanging factors
gives the exterior orientation sign.

For $v=\sum_jv_je_j$, put $s_j=\sum_{r\leq j}v_re_r$, $s_0=0$,
and define the projected oriented triangle chain
\[
 A_v=\sum_{j=1}^6[0,s_{j-1},s_j],\qquad
 \partial A_v=\sum_jv_j\gamma_j-\ell_v.
\]
Indeed the diagonal edges cancel, and the edge from $s_{j-1}$ to
$s_j$ is the translated coordinate loop with signed winding $v_j$.
For the columns $q_j=Qe_j$ of either actual unimodular marking matrix,
set
\[
\begin{split}
\widehat H_{Q,k}(e_{i_1}\wedge\cdots\wedge e_{i_k})
=\sum_{r=1}^k(-1)^{r-1}
 &\iota(q_{i_1})*\cdots*\iota(q_{i_{r-1}})*A_{q_{i_r}}\\
 &*\ell_{q_{i_{r+1}}}*\cdots*\ell_{q_{i_k}}.
\end{split}
\]
The Leibniz sign cancels the displayed sign. The sum telescopes to
the exact identity
\[
 \partial\widehat H_Q=\iota\bigwedge Q-Q_*\iota.
\]
Thus every lower-degree homotopy below is a specified finite sum
of integer triangles and products.

\item \emph{The prescribed top component and its composition law.}
Orient the original torus by $e_1\wedge\cdots\wedge e_6$ and put
$J(z)(u)=\operatorname{vol}(u\wedge z)$. The cooriented torus
$H_j=\{x_j=0\}$ has class
$(-1)^{j-1}e_1\wedge\cdots\widehat e_j\cdots\wedge e_6$.
Consequently
\[
 J\bigwedge^5Q=(\det Q)Q^{-t}J,\qquad
 \mathbf1^tJ=(-1,1,-1,1,-1,1).
\]
For a primitive covector $l$, the integer density
$\lfloor l^tx\rfloor$ on the fundamental cube defines $D_l$.
The jumps on internal cuts and opposite faces give
\[
 \partial D_l=\sum_jl_jH_j-H_l,
\]
where $H_l$ is the cooriented subtorus $l^tx=0\pmod1$.
Writing $\iota^\vee(c)=\sum_jc_jH_j$, define
\[
 D_Q(c)(x)=(\det Q)c^t\lfloor Q^{-1}x\rfloor.
\]
It is the integer chain $(\det Q)\sum_jc_jD_{Q^{-t}e_j}$, and
\[
 \partial D_Q(c)=\iota^\vee((\det Q)Q^{-t}c)
                         -Q_*\iota^\vee(c).
\]
The homotopy used from now on is specified in all degrees by
\[
 H_{Q,k}=\widehat H_{Q,k}\ (1\leq k\leq4),\quad
 H_{Q,5}(z)=D_Q(Jz),\quad H_{Q,0}=H_{Q,6}=0.
\]
It satisfies $\partial H_Q=\iota\bigwedge Q-Q_*\iota$.
In degree six both maps send the fundamental chain to $\det Q$
times that chain. Since $d_L=0$, this homotopy equation has no
adjacent-degree source term; the specified degree-five replacement
therefore preserves every lower-degree equation.

These floor chains satisfy the exact composition law
\[
 D_{PQ}(c)=D_P((\det Q)Q^{-t}c)+P_*D_Q(c).
\]
To verify it put $a=P^{-1}x$, $m=\lfloor a\rfloor$, $a_0=a-m$.
The integer identity
$\lfloor Q^{-1}a\rfloor=Q^{-1}m+\lfloor Q^{-1}a_0\rfloor$,
multiplied by $(\det P)(\det Q)c^t$, is precisely the displayed
identity of oriented densities. A common rational subdivision gives
the equality of integral chains, including orientation-reversing
$P,Q$. In particular $D_I=0$ and the inverse law is exact.

\item \emph{The actual finite-cap squares.}
Use the literal matrices of P1.12:
\[
 Q_i=P_i\operatorname{diag}(1,K_i),\quad
 \widetilde M_i=Q_i^{-1}M_iQ_i,\quad \det Q_i=-1.
\]
The first signed monodromy fixes the first coordinate and positively
cycles the other five. The second fixes the first, positively cycles
the next four, and negates the sixth. Their product-cell actions
are exactly the exterior matrices. In original coordinates put
\[
 \widetilde\iota_i=Q_{i*}\iota\bigwedge Q_i^{-1},\qquad
 M_{i*}\widetilde\iota_i=\widetilde\iota_iW_i.
\]

The geometric mapping-cylinder complex is
\[
 \mathcal B(M)_k=\mathcal P_k(T^6)\oplus\mathcal P_{k-1}(T^6),
 \qquad
 \partial(a,x)=(\partial a+(M_*-I)x,-\partial x).
\]
Its map to actual chains sends the second component to its oriented
edge prism. Cutting along the fibre identifies the relative chains
with the shifted fibre chains, so this map is an integral homology
equivalence. Use the same two-prism construction for $\mathcal G$.
The signed-product boundary comparison is
$\widetilde\Phi_i(f,x)=(\widetilde\iota_if,\widetilde\iota_ix)$.
Let $\jmath_i:E_i\to\mathcal P_*(E_i^{\rm geom})$ be the cap
product-cell inclusion, transported by $K_i$ to its original basis.

The actual collar map in the supplied marking is
\[
 q_i(t,y,\theta)=[y,\theta+\mu_it],\quad
 \mu_i=\varepsilon_im_i,\quad (m_1,m_2)=(5,4),\quad
 (\varepsilon_1,\varepsilon_2)=(1,-1).
\]
At $t=y=0$ it is the prescribed positive cap-base circuit
$[0,\theta]$. Its $t$ circle has degree $\mu_i$.
Subdividing the swept base interval proves on every cell that
\[
 q_i\widetilde\Phi_i=\jmath_i b_i,
\]
where $b_i$ is the literal map (1.12.5): a fibre cell in $K$ remains
that cell, the $t$ circle produces the signed norm of its $m_i$
strips, and the $\theta$ circle produces one strip. A cell containing
both directions has dimension-dropping image because they enter the
single sum $\theta+\mu_it$. This proves all four rows, including
the based affine meridian and the zero mixed row.

Define
\[
 H_i=H_{Q_i}\bigwedge Q_i^{-1},\quad
 \partial H_i=\iota-\widetilde\iota_i,\quad
 \mathcal K_i^Q=H_iW_i-M_{i*}H_i.
\]
Then $\partial\mathcal K_i^Q=\iota W_i-M_{i*}\iota$, and the
common-fibre comparisons are
\[
\begin{split}
 \Phi_i^Q(f,x)&=(\iota f+\mathcal K_i^Qx,\iota x),\\
 \Phi_G^Q(f,x_1,x_2)&=
 (\iota f+\mathcal K_1^Qx_1+\mathcal K_2^Qx_2,
                                      \iota x_1,\iota x_2).
\end{split}
\]
The full cylinder homotopy is
\[
 \mathcal T_i(f,x)=(H_if,-H_ix),\qquad
 \Phi_i^Q-\widetilde\Phi_i
       =\partial\mathcal T_i+\mathcal T_i d_{B_i}.
\]
Its fibre component is
$(\iota-\widetilde\iota_i)f+H_i(W_i-I)x-(M_{i*}-I)H_ix$,
and its edge component is $(\iota-\widetilde\iota_i)x$.
Applying the actual map gives the cap homotopy
\[
 q_i\Phi_i^Q-\jmath_i b_i=\partial Z_i+Z_i d_{B_i},
 \qquad Z_i=q_i\mathcal T_i.
\]
This specifies the square in every degree using the preceding
integer triangles, floor chains and actual collar map.

\item \emph{The integral cap-top-cell correction.}
For $\widetilde M_1$, the floor chain is zero. For
$\widetilde M_2$, its negative sixth coordinate gives
\[
 D_{\widetilde M_2}(c)=-c_6[T^6].
\]
In increasing exterior order $z=e_1\wedge\cdots\wedge e_5$
has $Jz=-e_6$. Its suspension has differential $-2z$, and its
floor-chain lift adds one fibre top cell. The composition law and
$M_iQ_i=Q_i\widetilde M_i$ give the exact transported equation
\[
 D_{M_i}J-\mathcal K_i^Q
   =Q_{i*}D_{\widetilde M_i}J\bigwedge^5Q_i^{-1}=h_i[T^6],
 \qquad h_1=0,\quad h_2=(1,-1,-1,1,-1,0).
\]
The signed source row is $(1,0,0,0,0,0)$; its raw transport is
$(\det Q_2)(1,0,0,0,0,0)\bigwedge^5Q_2^{-1}$, exactly $h_2$.
Let $U_i$ be identity outside degree six and
$U_{i,6}(f,x)=(f+h_ix,x)$. Define $\mathcal K_i$ to equal
$\mathcal K_i^Q$ outside source degree five and $D_{M_i}J$ there.
The new comparisons satisfy
\[
 \Phi_i=\Phi_i^QU_i,\qquad \Phi_G=\Phi_G^QU_G,
 \quad U_{G,6}(f,x_1,x_2)=(f+h_1x_1+h_2x_2,x_1,x_2).
\]

The actual top-cell map of the second cap has degree $\mu_2=-4$.
The source change therefore adds $-4[E_2]$. Since the changed
suspension has differential $-2z$, the integral homotopy is
$z\mapsto2[E_2]$. Transporting both source and target orientations
gives
\[
 \Gamma_1=0,\quad
 \Gamma_2=2h_2=(2,-2,-2,2,-2,0):L_5\longrightarrow E_{2,6},
 \qquad b_iU_i-b_i=d_{E_i}\Gamma_i+\Gamma_i d_{B_i}.
\]
Extend $\Gamma_i$ by zero on the other boundary summands. The
nonzero equation is $\Gamma_2(W_{2,5}-I)=-4h_2$, and
$d_E\Gamma_2=0$. The geometric homotopy value is the specified cap
fundamental chain. Because $H_{i,6}=0$,
$\mathcal T_i(U_i-I)=0$, so the complete finite-cap square is
\[
 q_i\Phi_i-\jmath_i b_i
  =\partial Z_i^{\rm pol}+Z_i^{\rm pol}d_{B_i},\qquad
 Z_i^{\rm pol}=Z_i+\jmath_i\Gamma_i.
\]
All changes are identity in degree one and retain the based meridian.

\item \emph{The simultaneous complement comparison.}
Take the numerical cone $Y=\operatorname{Cone}(j_G,-b_1,-b_2)$
of (1.12.6), and the geometric cone $\mathcal Y$ of the two actual
boundary maps to $\mathcal G$ and the two actual maps $q_i$.
The map $\Psi_Y:Y\to\mathcal Y$ is $\Phi_G$ on the graph,
$\jmath_i$ on the caps, $\Phi_i$ on the shifted boundaries, and
$+Z_i^{\rm pol}$ from a shifted boundary to its cap summand.
The sign follows from
$\Phi_DF-F^{\rm geom}\Phi_B=\partial Z+Zd_B$
for the cone differential $(d_D,F;0,-d_B)$.
The diagonal fibre, cylinder and cap maps are integral homology
equivalences; the relative cone sequence proves the same for
$\Psi_Y$. Prisms on the actual collars identify $\mathcal Y$ with
the chains of the actual complement by another such equivalence.

These are integral chain-homotopy equivalences as well. The complexes
are free. The cone of each displayed quasi-isomorphism is free and
acyclic; choose a basis of its boundary subgroups and integral
sections $s_k$ of the differential. In the decomposition
$P_k=B_k\oplus s_{k-1}B_{k-1}$, the map
$T(b+s b')=s_kb$ satisfies $\partial T+T\partial=I$.
For a free-homology contraction retain the specified cycle summand
$\iota H_k$ and project onto it; the same formula gives
$\partial T+T\partial=I-\iota p$ and $p\iota=I$.
The maps and diagram homotopies being inverted are already fixed
above; this contraction does not change their markings.

The numerical relation to the preceding common-fibre comparison is
\[
 \Psi_Y=\Psi_Y^Q S_Y.
\]
Here $S_Y$ is $U_G$ on the graph, identity on caps, $U_i$ on the
shifted boundaries, and $+\Gamma_i$ from shifted $B_{i,5}$ to
$E_{i,6}$. Its correction squares to zero, so its inverse subtracts
that correction. The exact literal matrices verify every
differential equation. Thus the actual finite-cap diagram has been
transported with its source maps, target maps and homotopies specified.

\item \emph{The actual cusp inverse word.}
For $g_0=g_2^{-1}g_1^{-1}$ extend the two chosen homotopies by
\[
 \mathcal K_{gh}=\mathcal K_gW_h+M_{g*}\mathcal K_h,\quad
 \mathcal K_{g^{-1}}=-M_{g^{-1}*}\mathcal K_gW_{g^{-1}},\quad
 \mathcal K_e=0.
\]
Free cancellation preserves these formulas, and
$\partial\mathcal K_g=\iota W_g-M_{g*}\iota$.
In source degree five the exact floor composition law gives
$\mathcal K_g=D_{M_g}J$. Define
\[
 \Phi_0(f,x)=(\iota f+\mathcal K_{g_0}x,\iota x).
\]
The actual inverse-word map on cylinder chains is
\[
 \mathcal J_0(a,x)=(a,-M_{2*}M_{0*}x,-M_{0*}x).
\]
This follows by subdividing the actual source prism into its two
inverse-edge strips, with the transports and zero clutch translation
of (1.12.3). Its precise comparison with that numerical map is
\[
\begin{split}
 \mathcal R(f,x)&=(0,\mathcal K_{g_1^{-1}}x,\mathcal K_{g_0}x),\\
 \Phi_Gj_{0,G}-\mathcal J_0\Phi_0
      &=\partial\mathcal R+\mathcal R d_{B_0}.
\end{split}
\]
For the edge components this is the equation for
$\partial\mathcal K_g$; the fibre component is
\[
 -\mathcal K_1W_2W_0-\mathcal K_2W_0-\mathcal K_{g_0}
 =(M_{1*}-I)\mathcal K_{g_1^{-1}}
                      +(M_{2*}-I)\mathcal K_{g_0},
\]
obtained directly from the word formulas. Also
$\mathcal R d_{B_0}=0$ because $\mathcal R$ vanishes on the fibre.
Graph inclusion carries this exact homotopy into the complement.

The numerical top-coordinate relation is
\[
 h_0=-h_1W_{2,5}W_{0,5}-h_2W_{0,5}=(-1,1,1,-3,-1,0),
 \quad U_{0,6}(f,x)=(f+h_0x,x),
\]
with identity in other degrees, and
\[
 \Phi_0=\Phi_0^QU_0,\qquad S_Yj_Y=j_YU_0.
\]
In particular $j_Y$ is transported explicitly. On the actual three
invariant five-vector columns $X$, $h_0X=(0,-2,-2)$; this nonzero
coordinate is retained by these maps.

\item \emph{The two cusp suspension values needed for the reduction.}
The invariant integral lattice in $L_5$ has basis
\[
 x_j=a_1\wedge a_2\wedge a_3\wedge
 b_1\wedge\cdots\widehat{b_j}\cdots\wedge b_3,
 \qquad j=1,2,3.
\]
This follows from the integral dual action of the shear;
$l_j=Jx_j$ are primitive covectors with $M_0^{-t}l_j=l_j$.
Their subtori are the marked $b_j=0$ tori. The floor densities give
the stronger exact identity
\[
 D_{M_0}(l_j)=(I-M_{0*})D_{l_j}.
\]
Indeed for $m=\lfloor M_0^{-1}x\rfloor$ the right side has density
\[
 \lfloor l_j^tx\rfloor-
 \lfloor l_j^t(M_0^{-1}x-m)\rfloor=l_j^tm.
\]
Consequently the comparison with the actual invariant-subtorus
mapping torus is the integral boundary
\[
 \Phi_0(0,x_j)-(0,H_{l_j})
 =((I-M_{0*})D_{l_j},\partial D_{l_j})
 =\partial(0,-D_{l_j}).
\]
This specifies the seven-chain and fixes the fibre coordinate.
The inverse-word homotopy just displayed tracks that same cycle
into the complement. Under (1.19.1) its cusp image has base
projection contained in the two-dimensional torus $b_j=0$ and
angular dimension at most three. Subdivision by the finite
moment-cell stratification gives an image of dimension at most five,
so its sixth homology class vanishes. The permitted deformation of
bounded period units transports this marked map by homotopy.

In degree one, $H_{Q,0}=0$ gives $\mathcal K_{g,0}=0$.
Thus $\Phi_0(0,1)$ is exactly the zero-section meridian; its
zero-section disc in the cusp cap proves its image is zero.
Both critical suspension values are now proved in the simultaneous
marked diagram, with the finite meridians fixed as above.

\item \emph{The remaining suspension values and their integral shears.}
By P1.11 the cusp homology is free, of ranks
$(1,3,9,9,10,3,1)$, and the collapse $\kappa_k$ has primitive
image, onto except for its rank-nine image in degree four.
Its source basis is $(a_1,a_2,a_3,b_1,b_2,b_3)$; the actual matrix
$S$ displayed in P1.1 transports it to the original marking by
\[
 \kappa_k^{\rm orig}=\kappa_k\bigwedge^kS^{-1}.
\]
Integral Smith reduction of $W_{0,k}-I$ gives free boundary ranks
$(1,4,12,18,18,12,4,1)$ and just one nonunit pair
$t\mapsto2z$, with $\deg t=4$, $\deg z=3$.
Choose the free fibre cycles as the quotient basis detected by
$\kappa$, free suspension cycles, and this pair, and contract
the unit pairs integrally. The cusp complex contracts onto its free
homology by the boundary/section formula above. Its actual
transferred map has fibre part $\kappa$ and definite suspension
values, denoted $a_k(s)$.

For the actual smooth eight-dimensional cusp cap $W$, local relative
duality gives
$H_j(W,B_0;\mathbb Z)=H^{8-j}(C;\mathbb Z)$, all free.
The surjection from the fibre to $H_5(C)$ makes the boundary map
surjective there. The pair sequence therefore injects the rank-nine
$H_5(W,B_0)$ into the rank-eighteen $H_4(B_0)$.
Its image in $H_4(C)$ has rank nine and contains the primitive
rank-nine collapse image, so the two images are equal. Thus
$a_k(s)\in\operatorname{im}\kappa_k$ for $2\leq k\leq5$;
the other three degrees follow from surjectivity of $\kappa_k$.

Let $J_k:L_k\to Y_k$ be the literal fibre inclusion. The saved
integer matrices $F_k$ satisfy
\[
 d_YF_k=J_k\qquad(k=2,3,4,5).
\]
They are obtained from the integral cycle and incoming Smith bases:
project each column of $J_k$ to the cycle basis, apply the incoming
left Smith matrix, divide its nonzero unit entries, and apply the
corresponding incoming right basis. The homology coordinates are
zero. Every resulting entry and the displayed equality are checked
in \path{code/current_audits/audit_p119_fresh/fresh_saved_verification.json}.
The map $\Psi_Y$ sends this literal $J_k$ to the actual common
fibre, so these are also the required geometric complement fillers.

Choose the integral section $\iota_k^\kappa$ of $\kappa_k$ on
its primitive image. Replace each free suspension cycle by
$s-\iota_k^\kappa a_k(s)$. This is an integral boundary automorphism
fixing the fibre, with inverse obtained by adding the same term.
The complement change is the specified homotopy
$-F_k\iota_k^\kappa a_k$; it is zero on fibre generators, so
its composition with the remaining differential is zero.
These operations remove the free suspension values in degrees
two through five. The values in degrees one and six have already
been proved zero. For a homotopy $F'-F=dT+Td$ the exact cone
isomorphism is $(a,b)\mapsto(a-Tb,b)$, with inverse adding $Tb$.
Use it also for the inverse-word homotopy $\mathcal R$ above.
The normalized homology maps are therefore those of
$f_0=(\kappa,0)$ in the full actual gluing diagram.

